(a) On an affine open $\operatorname{Spec}A=t(U)$, the residue field at $\mathfrak p$ is $\operatorname{Frac}(A/\mathfrak p)$. If $\mathfrak p$ is maximal, this is $k$ by the Nullstellensatz. If it is $k$ as a $k$-algebra, then $A/\mathfrak p$, lying between $k$ and $k$, is $k$, so $\mathfrak p$ is maximal. The corresponding point of $t(V)$ is the singleton closed subset of $V$, and hence is closed globally.

(b) The local homomorphism $\mathcal O_{Y,f(P)}\to\mathcal O_{X,P}$ induces a $k$-algebra inclusion $k(f(P))\to k(P)=k$. It restricts to the identity on $k$, so $k(f(P))=k$.

(c) A morphism of varieties determines its scheme morphism, whose restriction to closed points is the original morphism. Thus the natural map is injective. Conversely, a $k$-morphism $f:t(V)\to t(W)$ sends closed points to closed points by (a) and (b), giving $\varphi:V\to W$. The correspondence between open sets of a variety and its image under $t$ shows that $\varphi$ is continuous. For a regular function $s$ on an open subset of $W$, the section $f^\#s$ has value $s(\varphi(P))$ at every $P\in V$, because the induced residue-field maps are the identity on $k$. Hence $\varphi$ pulls regular functions back to regular functions and is a morphism of varieties.

The maps $f$ and $t(\varphi)$ have the same inverse images of opens, since opens of $t(V)$ are determined by their closed points. Their maps on sections agree, since regular functions are determined by their values at points of $V$. Thus $f=t(\varphi)$, proving surjectivity.
